% !TEX root = tarasov-qasp-zh.tex
% 章首扉图与格言（sec-05，守恒定律）在 zh/ch-09-zh.tex 末尾——照英文源物理位置（ch-09.tex 尾部有效代码）。
\chapter[你能运用能量守恒定律和动量守恒定律吗？]{你能运用能量守恒定律\\和动量守恒定律吗？\label{ch-10}}

\begin{dialogue}

\tchr 首先，我想提出几个简单的问题。
\begin{description}[leftmargin=1cm]
\item[第一个问题：] \hlred{若干物体无摩擦地沿两个高度同为 $H$、倾角分别为 $\alpha_{1}$ 与 $\alpha_{2}$ 的斜面滑下。物体的初速度为零。求物体到达斜面底端时的速度。}

\item[第二个问题：] \hlred{我们知道，用加速度和所通过的路程表示物体末速度的公式 $v = \sqrt{2as}$，只适用于初速度为零的情形。如果物体具有初速度 $v_{0}$，这个公式应是什么样子？}

\item[第三个问题：] \hlred{物体从高度 $H$ 处以水平速度 $v_{0}$ 抛出。求它到达地面时的速度。}

\item[第四个问题：] \hlred{物体以初速度 $v_{0}$、与水平方向成角 $\alpha$ 斜向上抛出。求它在飞行中达到的最大高度。}
\end{description}

\stda 我用下面的办法解第一个问题。先看两个斜面中的一个，比如倾角为 $\alpha_{1}$ 的那个。物体受两个力：重力 $P$ 和支承力 $N_{1}$。我们把力 $P$ 分解为两个分力：一个沿斜面方向（$P \sin \alpha_{1}$），另一个垂直于斜面（$P \cos \alpha_{1}$）。然后对垂直于斜面方向的力写出方程
\begin{equation*}%
P \cos \alpha_{1} - N_{1} = 0
\end{equation*}
对沿斜面方向的力写出方程
\begin{equation*}%
P \sin \alpha_{1}  = \frac{P a_{1}}{g}
\end{equation*}
其中 $a_{1}$ 是物体的加速度。由第二个方程得 $a_{1} =g \sin \alpha_{1}$。物体通过的路程为 $H/\sin \alpha_{1}$。再利用第二个问题中提到的公式，求得物体在路程末端的速度
\begin{equation*}%
v_{1}  = \sqrt{2a_{1}s_{1}} = \sqrt{\frac{2gH \sin \alpha_{1}}{\sin \alpha_{1}}} = \sqrt{2gH}
\end{equation*}
由于最终结果与倾角无关，它对倾角为 $\alpha_{2}$ 的第二个斜面同样适用。
\pnote{底本此式首项误排为 $v_{1}=\sqrt{a_{1}s_{1}}$，脱漏因子 2；按学生所引的公式 $v=\sqrt{2as}$ 及式中后续两步，应为 $\sqrt{2a_{1}s_{1}}$，中文版径改。本章正文另有若干拼写级排印小误（如 ``at an angle a'' 漏排 $\alpha$、``energy; work'' 分号误植、``law your this your of'' 衍文、``equals, the square'' 逗号误植、``You' are''、``mass.~$M$'' 断句误植、$v_{o}$ 误用字母 o、木块误作 ``box'' 等），均径改不注。}

为解第二个问题，我要利用熟知的运动学关系式
\begin{equation*}%
\begin{split}
v & = v_{0} + at\\
s & = v_{0}t+ \frac{at^{2}}{2}
\end{split}
\end{equation*}
由第一个方程得 $t= (v - v_{0})/a$。把它代入第二个方程，得
\begin{equation*}%
s = \frac{ v_{0}(v - v_{0})}{a} + \frac{a}{2} \frac{(v - v_{0})^{2}}{a^{2}}
\end{equation*}
即
\begin{equation*}%
2sa = 2 v_{0}v - 2v_{0}^{2}+ v^{2} - 2 v_{0}v + v_{0}^{2}
\end{equation*}

由此得 $2sa=v^{2}-v_{0}^{2}$。最终结果是
\begin{equation}%
v= \sqrt{2as+v_{0}^{2}}
\label{eq-48}
\end{equation}
解第三个问题时，我先求末速度的水平分量 $v_{1}$ 和竖直分量 $v_{2}$。由于物体沿水平方向做匀速运动，$v_{1} = v_{0}$。在竖直方向上，物体以加速度 $g$ 运动，但没有初速度，因此可以用公式 $v_{2}=\sqrt{2gH}$。由于直角三角形两条直角边的平方和等于斜边的平方，最终答案是
\begin{equation}
v=\sqrt{v_{1}^{2}+v_{2}^{2}} =\sqrt{v_{0}^{2}+2gH}
\label{eq-49}
\end{equation}


第四个问题已经在\chap{ch-05} 中讨论过了。需要把初速度分解为水平（$v_{0} \cos \alpha$）和竖直（$v_{0} \sin \alpha$）两个分量。然后考察物体的竖直方向运动：首先由匀减速运动中速度随时间变化的关系式（$v_{v}=v_{0} \sin \alpha -gt$）求上升时间 $t_{1}$，注意当 $t=t_{1}$ 时物体的竖直速度为零。于是 $v_{0} \sin \alpha -gt_{1} =0$，由此 $t_{1} = (v_{0}/g) \sin \alpha$。知道了 $t_{1}$，再由匀减速运动中路程随时间变化的公式，求出所达到的高度 $H$。于是
\begin{equation*}%
H = v_{0} t_{1} \sin \alpha - \frac{gt_{1}^{2}}{2} =  \frac{v_{0}^{2}}{2g}\sin^{2} \alpha
\end{equation*}
\pnote{底本此式末因子误排为 $\sin \alpha$，漏平方指数；按推导应为 $\sin^{2}\alpha$，且与教师随后用能量法所得结果一致（教师明言四个答案都对），中文版径改。}

\tchr 四个问题你都得到了正确的答案。不过，我对你解这些题的方式并不满意。如果运用能量守恒定律（law of conservation of energy），这些题都可以解得简单得多。你自己看吧。
\begin{description}[leftmargin=1cm, style=nextline]

\item[第一个问题。] 能量守恒定律的形式是 $mgH = mv^{2}/2$（物体在斜面顶端的势能等于它在斜面底端的动能）。由此方程我们立即求得物体到达底端时的速度
\begin{equation*}
v = \sqrt{2gH}
\end{equation*}


\item[第二个问题。] 在这种情形下，能量守恒定律的形式是
$mv_{0}^{2}/2+mas= mv^{2}/2$，其中 $mas$ 是使物体获得加速度 $a$ 的那些力所做的功。由此立即可得 $v_{0}^{2} + 2as = v^{2}$，即最终为
\begin{equation*}
v=\sqrt{2as+v_{0}^{2}}
\end{equation*}

\item[第三个问题。] 把能量守恒定律写成 $mgH+mv_{0}^{2}/2= mv^{2}/2$。于是结果为
\begin{equation*}
v=\sqrt{2gH+v_{0}^{2}}
\end{equation*}

\item[第四个问题。] 在抛出点，物体的能量等于 $mv_{0}^{2}/2$。在轨迹的最高点，物体的能量是 $mgH+mv_{1}^{2}/2$。由于最高点的速度 $v_{1}$ 等于 $v_{0} \cos \alpha$，利用能量守恒定律
\begin{equation*}
\frac{mv_{0}^{2}}{2} = mgH + \frac{mv_{0}^{2}}{2} \cos^{2} \alpha
\end{equation*}
求得
\begin{equation*}
H=\left(\frac{v_{0}^{2}}{2g} \right) \left( 1 - \cos^{2} \alpha \right)
\end{equation*}
即最终为
\begin{equation*}
H=\left(\frac{v_{0}^{2}}{2g} \right) \sin^{2} \alpha
\end{equation*}
\end{description}
\pnote{底本"第二个问题"一段的中间结果误排为 $v_{0}^{2}+2as=v_{2}$，应为 $v^{2}$（角标误作下标）；"第三个问题"一段的能量式误排为 $mv_{0}^{2}2$，脱除号，应为 $mv_{0}^{2}/2$。两处均按上下文径改。}

\stda 是的，这些题可以用简单得多的办法来解，这一点我很清楚。我只是没想到要用能量守恒定律。

\tchr 遗憾的是，考生们经常忘记这条定律。结果，他们开始用更繁琐的方法去解这类题，从而增加了出错的可能。我的忠告是：\hlblue{要更加灵活、更加广泛地运用能量守恒定律。}这就引出一个问题：你能多么熟练地运用这条定律？
\pnote{与其死记 $v^{2}-v_{0}^{2}=2as$ 这类运动学公式，不如认清它就是动能定理：合外力做的功等于动能的增量（人教版必修第二册第八章"动能和动能定理"一节）。教师下面用能量守恒定律重解第二个问题，用的正是这一思想。}

\stda 在我看来，这不需要什么特别的本领；能量守恒定律本身是相当简单的。

\tchr 能不能正确地运用一条物理定律，并不取决于这条定律是复杂还是简单。来看一个例子。\hlred{设物体在水平面内沿圆周做匀速运动。没有任何摩擦力作用。物体受到向心力的作用。问：物体转一周的过程中，这个力做的功是多少？}

\stda 功等于力与力所通过的路程的乘积。于是在本题中，它等于 $(mv^{2}/R) 2 \pi R = 2 \pi mv^{2}$，其中 $R$ 是圆周的半径，$m$ 和 $v$ 是物体的质量和速度。

\tchr 按照能量守恒定律，功是不可能无影无踪地消失的。那么，你算出的这份功到哪里去了？

\stda 它被用来使物体转动了。

\tchr 我不明白你的意思。说得更确切一些。

\stda 它使物体保持在圆周上。

\tchr 你的推理不对。要把物体保持在圆周上，根本不需要做任何功。

\stda 那我就不知道该怎样回答您的问题了。

\tchr 传给物体的能量，用物理学家的话说，可以分配到下面几条"通道"上：
\begin{enumerate*}[label=(\arabic*)]
\item 增加物体的动能；
\item 增加物体的势能；
\item 该物体对其他物体所做的功；
\item 摩擦产生的热。
\end{enumerate*}
这就是那条普遍原理，而并非所有考生都能足够清楚地理解它。现在来看向心力的功。物体以恒定速度运动，因此它的动能并不增加，第一条通道关闭。物体在水平面内运动，因而它的势能不变，第二条通道也关闭。该物体没有对其他物体做任何功，所以第三条通道关闭。最后，各种摩擦都已被排除，这就关上了第四条、也是最后一条通道。
\pnote{这份"能量去向清单"就是中学里用能量观点解题时的基本账目：前两条通道之内的相互转化，对应人教版必修第二册第八章的机械能守恒定律（只有重力或弹力做功时机械能守恒）；第三、四条通道是机械能的出口——对外做功与摩擦生热（转化为内能）。教师这里逐条排除通道得出"向心力不做功"，是能量守恒定律典型的"对账"式用法。}

\stda 这么说来，向心力的功就根本没有去处了，是吗？

\tchr 正如你看到的，没有。现在该你在这件事上表明立场了。要么你承认能量守恒定律不成立，那样一来一切麻烦都没有了；要么你以这条定律成立为出发点，那么……不过，还是请你设法找到摆脱困难的出路吧。

\stda 我想，剩下的结论只能是：\hlblue{向心力根本不做任何功。}

\tchr 这是一个十分合乎逻辑的结论。我要指出，它是能量守恒定律的直接推论。

\stdb 这一切都不错，可是功的公式该怎么办呢？

\tchr 除了力和力所通过的路程，这个公式里还应包含力与速度方向之间夹角的余弦：
$ A = F s \cos \alpha$。在本题中，$\cos \alpha = 0$。
\pnote{功的公式 $A=Fs\cos\alpha$ 即人教版必修第二册第八章"功与功率"一节中的 $W=Fl\cos\alpha$，其中 $\alpha$ 是力与位移方向之间的夹角。力与位移（速度）垂直时功为零——这正是匀速圆周运动中向心力不做功的原因。}

\stda 啊，对了，这个余弦我完全忘掉了。

\tchr 我想再举一个例子。考虑两个由一根带旋塞的细管连通的容器——连通器（communicating vessels）。设起初全部液体都在左边的容器里，液柱高度为 $H$（\figr{fig-43}\drkgry{(a)}）。然后打开旋塞，液体从左边流入右边的容器。当两个容器中的液面相平、各为 $H/2$ 时，流动停止（\figr{fig-43}\drkgry{(b)}）。

我们来计算液体在初位置和末位置的势能。为此，把每个容器中液体的重量乘以液柱高度的一半。在初位置，势能等于 $PH/2$；在末位置，它是
\begin{equation*}
 \left( \frac{P}{2} \right)  \left( \frac{H}{4} \right) +  \left( \frac{P}{2} \right)   \left( \frac{H}{4} \right) =   \left( \frac{PH}{4} \right).
\end{equation*}
\begin{figure}[!h]%
\centering
\includegraphics[width=0.4\textwidth]{figs/ch-10/fig-043a.pdf}
\caption[连通器中液柱势能的问题。]{连通器中液柱势能的问题。(a) 全部液体在左容器；(b) 两容器液面相平；(c) 全部液体流入右容器。\label{fig-43}}
\end{figure}

于是，在末状态，液体的势能只有初状态的一半。有一半能量消失到哪里去了？

\stda 我试着按您刚才的建议来推理。势能 $\dfrac{PH}{4}$ 可能消耗在：对其他物体做的功、摩擦产生的热，以及液体自身的动能。是这样吗？

\tchr 完全正确。继续。

\stda 在本题中，从一个容器流向另一个容器的液体不对其他物体做功。液体在末状态没有动能，因为它处于静止。那么，只能得出结论：有一半势能转化成了摩擦产生的热。只是，对这究竟是哪一种摩擦，我的概念并不很清楚。

\tchr 你推理得正确，结论也得对了。关于这种摩擦的本性，我想补充几句。可以设想液体分成一层一层的，每一层各有确定的流动快慢。\hlblue{液层离管的内壁越近，它的速度越低。}各层之间有分子交换，结果，定向运动速度较大的分子来到定向运动速度较小的分子中间，反过来也一样。于是，"较快"的层对"较慢"的层起加速作用；反过来，"较慢"的层对"较快"的层起减速作用。

这一图景使我们可以谈论各层之间存在一种特殊的内摩擦（internal friction）。管子中部与管壁附近各层的速度相差越大，这种内摩擦的作用越强。请注意，靠近管壁的各层的速度，还与液体分子同管壁分子之间相互作用的性质有关。如果液体浸润（wetting）管壁，那么直接紧贴管壁的那一层实际上是静止的。
\pnote{内摩擦即液体的粘性，中学阶段不展开，大学普通物理中作为粘性流体处理；"分层流动、靠壁一层速度最低"的图像即层流的速度分布。浸润与不浸润现象在人教版选择性必修第三册第二章"液体"一节讨论（水浸润玻璃，水银不浸润玻璃）；"紧贴管壁的液层静止"在流体力学中称为无滑移条件。}

\stda 这是不是说，末状态液体的温度应当比初状态略高一些？

\tchr 是的，正是如此。现在我们把问题的条件稍加改变。设液体与管壁之间没有相互作用。于是，所有液层将以同一速度流动，内摩擦也就不复存在。这时，液体将怎样从一个容器流向另一个容器呢？

\stda 这时，势能将因液体获得动能而减小。换句话说，\figr{fig-43}\drkgry{(b)} 所示的状态并不是静止状态。液体将继续从左容器流向右容器，直到达到\figr{fig-43}\drkgry{(c)} 所示的状态。在这个状态里，势能重新与初始状态\figr{fig-43}\drkgry{(a)} 相同。

\tchr 在这之后液体又会怎样呢？

\stda 液体将开始倒流，从右容器流回左容器。结果，连通器中的液面将来回振荡。

\tchr 这样的振荡，例如可以在盛有水银的玻璃连通器中观察到。我们知道，水银不浸润玻璃。当然，振荡会随时间逐渐衰减，因为液体分子与管壁分子之间的相互作用不可能完全排除。

\stda 我明白了，能量守恒定律可以这样主动地加以运用。

\tchr 再给你们出一道题。\hlred{质量为 $m$ 的子弹沿水平方向以速度 $v_{0}$ 飞行，打中挂在绳上的质量为 $M$ 的木块，并陷在木块里。子弹打中之后，由于绳偏离平衡位置，木块将上升到多大的高度 $H$\drkgry{（\figr{fig-44}）}？}
\begin{figure}[!ht]
\centering
\includegraphics[width=0.4\textwidth]{figs/ch-10/fig-044a.pdf}
\caption[子弹打中悬挂在绳上的木块并陷入其中。]{子弹打中悬挂在绳上的木块并陷入其中。问题是求碰后木块上升的高度。\label{fig-44}}
\end{figure}

\stda 我们用 $v_{1}$ 表示子弹打中木块之后，木块带着子弹立即具有的速度。为求这个速度，我们利用能量守恒定律。于是
\begin{equation}%
\frac{m v_{0}^{2}}{2} = \left( m + M \right) \frac{v_{1}^{2}}{2}
\label{eq-50}
\end{equation}
由此得
\begin{equation}%
v_{1} = v_{0} \sqrt{\frac{m}{ m + M}}
\label{eq-51}
\end{equation}
知道了这个速度，我们再次借助能量守恒定律，求所要求的高度 $H$：
\begin{equation}%
(m+ M) gH = (m + M) \frac{v_{1}^{2}}{2}
\label{eq-52}
\end{equation}
式~\eqref{eq-50} 与式~\eqref{eq-52} 可以合并为
\begin{equation*}%
(m+ M) gH = \frac{mv_{0}^{2}}{2}
\end{equation*}
由此得
\begin{equation}%
H = \frac{m}{(m + M)} \frac{v_{0}^{2}}{2g}
\label{eq-53}
\end{equation}
\tchr （问学生乙）你觉得这个解法怎么样？

\stdb 我不同意这个解法。我们以前学过，这类情形应当用动量守恒定律（law of conservation of momentum）。因此，我会用另一个关系式代替式~\eqref{eq-50}：
\begin{equation}%
mv_{0} = (m+M) v_{1}
\label{eq-54}
\end{equation}
（子弹打中木块之前的动量，等于打中之后子弹和木块的动量）。由此得出
\begin{equation}%
v_{1} = v_{0} \frac{m}{(m + M)}
\label{eq-55}
\end{equation}
再利用能量守恒定律式~\eqref{eq-52}，把式~\eqref{eq-55} 的结果代入\eqref{eq-52}，得
\begin{equation}%
H = \left(\frac{m}{m + M} \right)^{2} \frac{v_{0}^{2}}{2g}
\label{eq-56}
\end{equation}
\tchr 我们得到了两种不同的意见、两个结果。问题在于：一个情形是子弹打中木块时运用了动能守恒定律，另一个情形运用的则是动量守恒定律。哪一个对？（问学生甲）你能说点什么来为自己的立场辩护吗？

\stda 我没想到要用动量守恒定律。

\tchr （问学生乙）那你怎么说？

\stdb 我不知道怎样论证我的立场。我只记得：\hlblue{处理碰撞问题时，动量守恒定律总是成立的，而能量守恒定律并不总成立。}既然在本题中这两条定律给出不同的结果，那么显然我的解法是对的。
\pnote{学生乙这句话中的"能量守恒定律"应理解为"机械能（动能）守恒"：能量守恒定律本身是普遍定律，永不失效。教师下文正是这样澄清的——把热 $Q$ 计入之后，总能量仍然严格守恒。人教版选择性必修第一册第一章"弹性碰撞和非弹性碰撞"一节的表述即如此：弹性碰撞中动量与机械能都守恒，非弹性碰撞中动量守恒而机械能有损失。}

\tchr 事实上，这个解法确实完全正确。不过，有必要把这件事看得更透彻一些。碰撞之后，相碰撞的物体连在一起（或一个进入另一个内部）一同运动，这样的碰撞称为"完全非弹性碰撞"（completely inelastic collision）。这类碰撞的特点是，相碰撞的物体中留下了永久形变（permanent set），结果放出一定的热量。因此，只涉及物体动能的式~\eqref{eq-50} 不适用。在本题中，必须用动量守恒定律式~\eqref{eq-54} 来求碰后木块带子弹的速度。
\pnote{这道题就是经典的"冲击摆"：碰撞阶段（子弹进入木块）用动量守恒，摆起阶段用机械能守恒——复合过程分阶段各用一条守恒定律，是这类问题的标准解法，也是高考与竞赛中的保留题型。学生甲的错，错在把能量守恒用到了碰撞阶段（该阶段机械能不守恒），而不是能量守恒定律本身失效。}

\stda 您是想说，能量守恒定律对完全非弹性碰撞不成立？可是这条定律是普遍适用的呀。

\tchr 毫无疑问，能量守恒定律对完全非弹性碰撞同样成立。碰后不守恒的是动能。我指的恰恰是动能，而不是全部能量。用 $Q$ 表示碰撞中放出的热，我们可以对上面讨论的完全非弹性碰撞写出如下的守恒定律方程组
\begin{equation}%
\left.
\begin{aligned}
mv_{0}  & = (m + M) v_{1} \\
\frac{mv_{0}^{2}}{2} & = \frac{(m+M) v_{1}^{2}}{2} + Q
\label{eq-57}
\end{aligned}
\right\}
\end{equation}
这里第一个方程是动量守恒定律，第二个是能量守恒定律（不仅包括机械能，还包括热）。方程组~\eqref{eq-57} 含有两个未知量：$v_{1}$ 和 $Q$。由第一个方程确定 $v_{1}$ 之后，就可以用第二个方程求放出的热 $Q$：
\begin{equation}%
Q = \frac{mv_{0}^{2}}{2} - \frac{(m+M)m^{2}v_{0}^{2}}{2(m+M)^{2}} =  \frac{mv_{0}^{2}}{2} \left( 1 - \frac{m}{m+M}\right)
\label{eq-58}
\end{equation}
由这个式子可以看出：质量 $M$ 越大，转化为热的能量越多。在极限情形，当质量 $M$ 无限大时，我们得到 $Q=m v_{0}^{2}/2$，即全部动能都转化为热。这十分自然：设想子弹打进一堵墙里。
\pnote{式~\eqref{eq-58} 的结果可整理为 $Q=\frac{1}{2}\,\frac{mM}{m+M}\,v_{0}^{2}$，其中 $\frac{mM}{m+M}$ 称为折合质量（reduced mass）——完全非弹性碰撞中损失的动能，等于折合质量的初动能。这是竞赛中常用的紧凑形式；还可以证明：在动量守恒的前提下，碰后共速时动能损失最大，所以完全非弹性碰撞是动能损失最大的碰撞。}

\stda 有没有不放出热的碰撞？

\tchr 有，这样的碰撞是可能的。它们称为"完全弹性碰撞"（perfectly elastic collision）。例如，两个钢球的碰撞可以以相当好的近似看作完全弹性碰撞。这时发生的是纯弹性形变，不放出热。碰撞之后，球恢复原来的形状。

\stda 您是说，在完全弹性碰撞中，能量守恒定律变成了动能守恒定律？

\tchr 是的，当然。

\stda 可是在这种情形下，我就不明白，怎样才能把动量守恒定律和能量守恒定律调和一致。对碰后的速度，我们会得到两个完全不同的方程。要不然，难道是动量守恒定律对完全弹性碰撞不成立？

\tchr 对完全弹性碰撞，两条守恒定律都成立：一条关于动量，一条关于动能。你们没有必要为这两条定律的调和担心，因为完全弹性碰撞之后，物体是以不同的速度飞开的。\hlblue{完全非弹性碰撞后，相碰撞的物体以同一速度运动（因为它们粘在一起）；而弹性碰撞后，每个物体各以自己确定的速度运动。}两个未知量需要两个方程。

我们来考虑一个例子。设质量为 $m$ 的物体以速度 $v_{0}$ 与静止的质量为 $M$ 的物体发生弹性碰撞。再设碰撞的结果是入射物体被弹回来。我们用 $v_{1}$ 表示碰后物体 $m$ 的速度，用 $v_{2}$ 表示物体 $M$ 的速度。于是动量守恒定律和能量守恒定律可以写成
\begin{equation}%
\left.
\begin{aligned}
mv_{0} & =Mv_{2}-  mv_{1} \\
\frac{mv_{0}^{2}}{2} & = \frac{Mv_{2}^{2}}{2} + \frac{mv_{1}^{2}}{2}
\label{eq-59}
\end{aligned}
\right\}
\end{equation}
注意第一个方程中的负号，它源于我们关于入射物体被弹回的假定。
\pnote{底本式~\eqref{eq-59} 第二个方程右端第一项误排为 $\frac{mv_{2}^{2}}{2}$。速度 $v_{2}$ 属于质量为 $M$ 的物体，该项应为 $\frac{Mv_{2}^{2}}{2}$；学生乙后文为斜面体问题写出的式~\eqref{eq-64} 中正是 $Mv_{2}^{2}/2$，可为佐证。中文版径改。}

\stdb 可是，碰撞后物体将沿哪个方向运动，您不可能总是事先知道。难道物体 $m$ 碰后不能以较小的速度沿原方向继续运动吗？

\tchr 这完全可能。在这种情形，解方程组~\eqref{eq-59} 时我们将得到负的速度 $v_{1}$。

\stdb 我想，物体 $m$ 碰后运动的方向由质量 $m$ 与 $M$ 之比决定。

\tchr 正是如此。若 $m<M$，物体 $m$ 将被弹回；若 $m=M$，它碰后将静止；若 $m>M$，它将以较小的速度沿原方向继续运动。不过在一般情形下，你不必为运动方向操心：只要假定某个方向并开始计算就够了。\hlblue{答案的符号自会指出你的错误（如果有错误的话）。}
\pnote{这三种情形可合并为一维弹性碰撞的通解：$v_{1}'=\frac{m-M}{m+M}v_{0}$，$v_{2}'=\frac{2m}{m+M}v_{0}$（人教版选择性必修第一册第一章"弹性碰撞和非弹性碰撞"一节有推导）。$m=M$ 时两球交换速度——牛顿摆演示的正是这一点。}

\stdb 我们知道，碰撞后两球可能互成一定角度飞开。而我们刚才假定的是沿同一条直线的运动。显然，那只是一种特殊情形。

\tchr 你说得对。我们讨论的是所谓对心碰撞（central collision）：球在碰撞前后都沿一条通过它们中心的直线运动。更一般的非对心碰撞（off-centre collision）的情形以后再讨论。眼下我想知道，是否一切都十分清楚了。
\pnote{对心碰撞在中文文献中习称"正碰"，非对心碰撞习称"斜碰"。人教版选择性必修第一册第一章讨论的主要是一维对心碰撞；二维斜碰问题按本节的方法，取两个相互垂直的方向分别写出动量守恒即可，如下面的例子。}

\stda 我想我现在明白了。依我看，在任何碰撞（弹性的或非弹性的）中，都有两条守恒定律可用：动量守恒和能量守恒。只不过是碰撞的性质不同，描述守恒定律的方程也就不同。\hlblue{处理非弹性碰撞时，必须计入碰撞放出的热。}

\tchr 你的话正确而中肯。

\stdb 据我理解，完全弹性碰撞和完全非弹性碰撞是两个极端情形。它们总适合于描述真实的情形吗？

\tchr 你提出这个问题是对的。我们考虑过的几种碰撞情形都是极端情形。真实的碰撞中总要多多少少产生一些热（理想的弹性形变并不存在），而且相碰撞的物体可能以不同的速度飞开。不过，在许多情形下，真实的碰撞用简化模型——完全弹性碰撞和完全非弹性碰撞——描述得相当好。
\pnote{真实碰撞介于两个极端之间。竞赛与大学教材中用恢复系数（coefficient of restitution）$e$ 描述：$e=1$ 为完全弹性碰撞，$e=0$ 为完全非弹性碰撞；$0<e<1$ 时，碰后的分离速度与碰前的接近速度之比为 $e$（牛顿碰撞公式）。}

现在来看非对心弹性碰撞的一个例子。\hlred{一个倾角为 \ang{45}、形状如斜面的物体，静置于水平面上。质量为 $m$ 的小球沿水平方向以速度 $v_{0}$ 飞行，与质量为 $M$ 的这个物体（斜面体）相碰。碰撞的结果，小球竖直向上弹起，而物体 $M$ 开始沿水平面无摩擦地滑动。求碰后小球开始竖直运动时的速度。}\drkgry{（\figr{fig-45}）。}

你们谁想来试试这道题？

\begin{figure}
\centering
\includegraphics[width=0.4\textwidth]{figs/ch-10/fig-045a.pdf}
\caption[小球与静止的斜面体相碰。]{小球与静止的、形状如斜面的物体相碰。问题是求碰后小球的速度。\label{fig-45}}
\end{figure}

\stdb 请允许我来。用 $v_{1}$ 表示所求的小球速度，用 $v_{2}$ 表示物体 $M$ 的速度。既然是弹性碰撞，我有权认为动能守恒。于是
\begin{equation}% 
%\left.
%\begin{split}
%mv_{0} & =Mv_{2}-  mv_{1} \\
\frac{mv_{0}^{2}}{2} = \frac{mv_{1}^{2}}{2} + \frac{Mv_{2}^{2}}{2}
\label{eq-60}
%\end{split}
%\right\}
\end{equation}
还需要一个方程，为此显然应当用动量守恒定律。我把它写成
\begin{equation}%
%\left.
%\begin{split}
mv_{0} =Mv_{2} +  mv_{1}
\label{eq-61}
%\end{split}
%\right\}
\end{equation}
不过，说实话，对最后一个方程我没有多大把握，因为速度 $v_{1}$ 与速度 $v_{2}$ 互相垂直。

\tchr 式~\eqref{eq-60} 是正确的。式~\eqref{eq-61} 不正确，正如你自己所预感的。应当记住，动量守恒定律是矢量方程，因为动量是矢量，具有与速度相同的方向。的确，当所有速度都沿同一条直线时，矢量方程可以换成标量方程——我们讨论对心碰撞时遇到的正是这种情形。\hlblue{但在一般情形下，必须把所有速度沿相互垂直的方向分解，并对其中每一个方向分别写出动量守恒定律}（如果问题在平面内考虑，矢量方程可以换成两个标量方程，即动量在两个相互垂直方向上投影的方程）。

对本题，我们可以取水平方向和竖直方向。沿水平方向，动量守恒定律的形式是
\begin{equation}%
mv_{0}= Mv_{2}
\label{eq-62}
\end{equation}
由式~\eqref{eq-60} 与式~\eqref{eq-62}，我们求得速度
\begin{equation*}%
v_{1} = v_{0}\sqrt{\frac{M-m}{M}}
%\label{eq-62}
\end{equation*}
\pnote{底本式~\eqref{eq-60} 右端第二项同样误排为 $\frac{mv_{2}^{2}}{2}$，应为 $\frac{Mv_{2}^{2}}{2}$（理由见式~\eqref{eq-59} 处评注）。又，底本由此解得的结果误排为 $v_{1}=v_{0}\sqrt{\frac{M-m}{m}}$：将式~\eqref{eq-60} 与式~\eqref{eq-62} 联立，解得 $v_{1}=v_{0}\sqrt{\frac{M-m}{M}}$；且当 $M\gg m$ 时 $v_{1}\to v_{0}$（相当于小球撞上固定的 \ang{45} 斜面，速度方向转了直角而大小不变），若分母为 $m$，则 $v_{1}$ 将随 $M$ 无限增大，物理上不可能。两处均径改，待考原书。}

\stdb 那么竖直方向怎么办呢？

\tchr 乍看起来，竖直方向的动量守恒定律似乎不成立。实际上它是成立的。碰前没有竖直方向的速度；碰后出现了一个竖直向上的动量 $mv_{1}$。我们很容易看出，问题中还有一个物体在参与：地球。如果没有地球，物体 $M$ 碰后就不会沿水平方向运动。用 $M_{\Earth}$ 表示地球的质量，用 $v_{\Earth}$ 表示地球因碰撞而获得的速度。

没有摩擦，这使我们可以把物体 $M$ 与地球的相互作用看作只发生在竖直方向。换句话说，地球的速度 $v_{\Earth}$ 竖直向下。于是，地球参与本题并不改变式~\eqref{eq-62} 的形式，而是引出一个描述竖直方向动量守恒定律的方程
\begin{equation}%
%\left.
%\begin{split}
mv_{1}  - M_{\Earth}v_{\Earth} = 0
\label{eq-63}
%\end{split}
%\right\}
\end{equation}
\stdb 既然地球也参与这个问题，那么看来能量关系式~\eqref{eq-60} 也必须修正了。

\tchr 那你打算对式~\eqref{eq-60} 作怎样的修正呢？

\stdb 我想补上一项，反映碰后地球的运动：
\begin{equation}%
%\left.
%\begin{split}
%mv_{0} & =Mv_{2}-  mv_{1} \\
\frac{mv_{0}^{2}}{2} = \frac{mv_{1}^{2}}{2} + \frac{Mv_{2}^{2}}{2} +  \frac{M_{\Earth}v_{\Earth}^{2}}{2}
\label{eq-64}
%\end{split}
%\right\}
\end{equation}

\tchr 你的想法完全合乎逻辑。不过，没有必要改变式~(\ref{eq-60})。事实上，由式~(\ref{eq-63}) 可知，地球的速度是
\begin{equation*}
v_{\Earth} = \frac{mv_{1}}{M_{\Earth}}
\end{equation*}
由于质量 $M_{\Earth}$ 实际上是无限大的，地球碰后的速度 $v_{\Earth}$ 实际上等于零。现在，我们把式~\eqref{eq-64} 中的项 $M_{\Earth}v_{\Earth}^{2}/2$ 改写成 $(M_{\Earth}v_{\Earth})v_{\Earth}/2$ 的形式。根据式~(\ref{eq-63})，乘积中的量 $M_{\Earth}v_{\Earth}$ 具有有限的值。把这个值乘以零（在本题中即乘以 $v_{\Earth}$），乘积也等于零。

由此我们可以得出结论：地球在本题中的参与方式非常特别。\hlblue{它获得了一份动量，但同时几乎没有得到能量。}换句话说，它参与动量守恒定律，却不参与能量守恒定律。这一情况特别醒目地表明：\hlblue{能量守恒定律和动量守恒定律是本质不同的两条定律，彼此独立。}
\pnote{严格地说，地球获得的动能并非零，只是与 $mv_{0}^{2}/2$ 相比小了约 $m/M_{\Earth}$ 倍，完全可以忽略。同样的图像适用于"皮球落到地面又弹起""蹬地起跳"等情形：地球吸收动量而几乎不吸收动能——这正是这些问题中机械能近似守恒、也无须担心"地球被撞动"的原因。}
\end{dialogue}

\cleardoublepage

\thispagestyle{empty}
\vspace*{-1.5cm}
\begin{figure*}[!ht]
\centering
\includegraphics[width=0.73\textwidth]{figs/sec/tarasov-qasp-sec-06.pdf}
\end{figure*}
\begin{fullwidth}
{\textsf{\Large \hlred{我们周围的世界充满了振动与波。在学习物理学中专门研究这些现象的分支时，请记住这一点。我们来讨论简谐运动，并以数学摆的振动作为特例；我们还将分析摆在非惯性参考系中的行为。
}}}
\end{fullwidth}

% NEWTERMS: law of conservation of energy | 能量守恒定律 | 首现于教师评述四题解法处
% NEWTERMS: law of conservation of momentum | 动量守恒定律 | 首现于学生乙评冲击摆解法处
% NEWTERMS: momentum | 动量 | 
% NEWTERMS: communicating vessels | 连通器 | 
% NEWTERMS: internal friction | 内摩擦 | 即粘性
% NEWTERMS: wetting | 浸润 | 底本用动词 wets/wet；人教版术语
% NEWTERMS: ordered motion | 定向运动 | 分子热运动的对偶概念
% NEWTERMS: completely inelastic collision | 完全非弹性碰撞 | 
% NEWTERMS: perfectly elastic collision | 完全弹性碰撞 | 底本亦作 completely elastic
% NEWTERMS: permanent set | 永久形变 | 
% NEWTERMS: central collision | 对心碰撞 | 习称正碰
% NEWTERMS: off-centre collision | 非对心碰撞 | 习称斜碰
% NEWTERMS: efficiency | 机械效率 | 习题 3
% NEWTERMS: ballistic pendulum | 冲击摆 | 仅评注用
% NEWTERMS: reduced mass | 折合质量 | 仅评注用
% NEWTERMS: coefficient of restitution | 恢复系数 | 仅评注用
